\section{Problem and evidence levels}
\label{sec:problem}

\subsection{Attention-defined saliency}
At layer $\ell$ and decode step $t$, let $\mathcal B_{\ell,t}$ contain the logical
KV blocks currently under consideration. A block groups contiguous original token
positions. We aggregate attention by averaging over heads and over a block's token
positions; denote the resulting statistic by $a_{\ell,b,t}$. It is a block-mean
attention statistic, not necessarily a mass normalized to sum to one over blocks.
For a horizon $h$ and a fraction $r$, the target is
$z_{\ell,b,t}^{(h)}=\mathbf 1[b\in\operatorname{Top}_{\lceil r|\mathcal B_{\ell,t}|\rceil}\{a_{\ell,b',t+h}:b'\in\mathcal B_{\ell,t}\}]$:
ranking is restricted to blocks present at $t$. We use $r=0.1$ and horizons
$h\in\{1,4,16,64\}$ (the archived collector's terminal-horizon caveat is in
\S\ref{sec:limits}).

A scorer $f(\mathbf x_{\ell,b,t})$ estimates the label probability from recorded
features; examples arrive in decode order, positives are a small fraction of blocks,
drift is a hypothesis to test, and serial dependence means millions of block rows are
not millions of independent observations.

\subsection{Three different evaluation targets}
Fig.~\ref{fig:paths} locates the evaluation paths.
Offline prediction scores future-attention labels on fixed full-cache traces;
served-oracle diagnostics score a current-attention target along the evaluated
trajectory, and attention importance is not answer correctness; closed-loop task
quality scores answers the selector's own decisions helped produce, here through
masking rather than a cache manager. A fraction of blocks allowed to attend is a
\emph{logical retention budget}, not a fraction of device memory released.

\begin{figure}[t]
\centering
\includegraphics[width=\linewidth]{evaluation_paths.pdf}
\caption{Evaluation paths. Runtime features and selector semantics differ
across paths (Table~\ref{tab:protocol}); physical-reference parity is tested separately.}
\label{fig:paths}
\end{figure}

\subsection{Statistical protocol}
Offline ranking uses AUC, threshold-based AUPRC (tied scores form one group, so
the value cannot depend on row order), precision/recall at $k=r$, Brier score and
expected calibration error (ECE). Top-$k$ metrics are reported pooled over test
rows \emph{and} per cache decision (request, layer, step) at the label's own
$\lceil r\,n\rceil$. Training and testing use disjoint requests and paired
bootstrap intervals resample requests, not block rows. Unless marked otherwise,
intervals are 95\%; equivalence intervals are 90\%. A 10,000-draw sensitivity
analysis jointly resamples all held-out model-requests sharing a source prompt,
on both the original request split and the source-disjoint split.
Top-$k$ score ties use a stable descending sort: pooled ties follow seeded sample
row order, and within-decision ties follow increasing block index. V2 label and
cross-layer indicator boundaries retain the collector's NumPy
\texttt{argpartition} tie rule (implementation-dependent on the original block
order), with exactly $\lceil rn\rceil$ positives; these are not randomized ties.
The main downstream
comparison aggregates paired request scores into architecture--dataset cells and a
TOST analysis~\cite{schuirmann1987tost,lakens2017equivalence} asks whether their
mean difference lies inside a pre-specified margin. Failure to detect a difference
is not equivalence, and equivalence of two comparisons with the same baseline does
not imply equivalence of the remaining pair.
